%%%%%%%%%%%%%%%%%%%%%%%%%%%%%%%%%%%%%%%%%
% Medium Length Professional CV
% LaTeX Template
% Version 2.0 (8/5/13)
%
% This template has been downloaded from:
% http://www.LaTeXTemplates.com
%
% Original author:
% Trey Hunner (http://www.treyhunner.com/)
%
% Important note:
% This template requires the resume.cls file to be in the same directory as the
% .tex file. The resume.cls file provides the resume style used for structuring the
% document.
%
%%%%%%%%%%%%%%%%%%%%%%%%%%%%%%%%%%%%%%%%%

%----------------------------------------------------------------------------------------
%	PACKAGES AND OTHER DOCUMENT CONFIGURATIONS
%----------------------------------------------------------------------------------------

\documentclass{resume} % Use the custom resume.cls style

%\usepackage[left=0.4 in,top=0.4in,right=0.4 in,bottom=0.4in]{geometry} % Document margins
\usepackage{zh_CN-Adobefonts_external} 
\usepackage{linespacing_fix}
\usepackage{cite}
\newcommand{\tab}[1]{\hspace{.2667\textwidth}\rlap{#1}} 
\newcommand{\itab}[1]{\hspace{0em}\rlap{#1}}

\name{Qian Xu} % Your name
%\address{123 Pleasant Lane \\ City, State 12345} % Your secondary addess (optional)
\address{+86 15600616741 \\ drqianxu@163.com}  % Your phone number and email

\begin{document}

%----------------------------------------------------------------------------------------
%	OBJECTIVE
%----------------------------------------------------------------------------------------

\begin{rSection}{Summary}

{目前在Intel担任机器学习工程师，主要着眼于Intel第三代VPU以及后续VPU的模型压缩技术研究（量化，稀疏化，剪枝等）。负责第三代VPU在0.8 MLperf inference的提交，同时为后续的VPU架构设计提供技术支持}


\end{rSection}
%----------------------------------------------------------------------------------------
%	EDUCATION SECTION
%----------------------------------------------------------------------------------------

\begin{rSection}{Education}

{\bf Phd of Electronics Engineering} \hfill {2013 - 2018}
\\ 
University of Chinese Academy of Sciences. Outstanding Graduates

{\textbf{Bachelor of Electronics Engineering}}  \hfill 2009 - 2013\\
Southeast University. Ranking 14/200+ 


%Minor in Linguistics \smallskip \\
%Member of Eta Kappa Nu \\
%Member of Upsilon Pi Epsilon \\


\end{rSection}
%----------------------------------------------------------------------------------------
%	TECHNICAL STRENGTHS SECTION
%----------------------------------------------------------------------------------------

\begin{rSection}{SKILLS}

\begin{tabular}{ @{} >{\bfseries}l @{\hspace{6ex}} l }
Skills &  Python, C++, C.  \\
Frameworks & Tensorflow, Pytorch
\end{tabular}

\end{rSection}


%----------------------------------------------------------------------------------------
%	WORK EXPERIENCE SECTION
%----------------------------------------------------------------------------------------

\begin{rSection}{PROJECTS}

\begin{rSubsection}{Provide compressed models for Intel’s 3rd generation VPU ‘Keembay’,}{September 2018 - Ongoing}{Team Leader}{}
\item Responsible for developing compression techniques to better utilize the inter-layer Mixed precision quantization (2, 4, 8) and sparisification ability of Keembay. 
\\ AutoML based model sparsification is used for weights and Activation Map compression is used for activations. PACT based method is used for activation quantization.
\item Responsible for providing ultra compressed models for MLPerf inference. Including Resnet-50, Mobilenet, SSD-Mobilenet, SSD-Resnet34, Bert and so on. 
\\Within 1\% top-1 accuracy drop of Resnet-50, we are able to achieve 84\% weight sparsity and 65\% activation sparsity, resulting in \textgreater 2$\times$ acceleration ratio compared to dense model.

\end{rSubsection} 


%------------------------------------------------

\begin{rSubsection}{Research on cutting edge model compression for next generation VPU design,}{January 2019 -  Ongoing}{Team Leader}{}
\item Implementing state of the art model quantization techniques, including low-bit quantization methods like PACT and SAWB,  Mixed precision methods including HAQ and HAWQ, Unlinear quantization methods like log-based, outlier-aware and Squantizer. 
\item Implementing state of the art model sparsification methods, including sparse momentum, activation map compression, auto model compression, lottery ticket and so on.
\item Provide innovative ideas to better compress the models. I come up with channelwise Hessian aware trace weighted quantization method, which outperforms state of the art methods by around 1\% top-1 accuracy.

\end{rSubsection}


\begin{rSubsection}{Profiling for next generation VPU,}{January 2019 -  Ongoing}{Participant}{}
\item Providing performance analysis for 3rd and next generation devices, including latency modeling for sparsification and mixed precision quantization.
\item Providing precision reference for hardware, including fp8 accuracy, winograd accuracy and so on.
\end{rSubsection}

\begin{rSubsection}{DNN based Epoxy defect detection,}{May 2020 -  Ongoing}{Team Leader}{}
\item Experiments on cutting edge object detection networks for defect detection rather than traditional image processing methods.
\item Use label smoothing and other metrics to fight against noisy label in the training set.
\item Image augmentation using traditional image processing metrics.
\end{rSubsection}

\begin{rSubsection}{Machine learning based Great Wall virtual repairing,}{June 2018 - December 2018}{Participant}{}
\item Cooperating with the Chinese Foundation for Cultural Heritage Conservation to launch an AI based method for its restoration.
\item Use 3D-GAN to obtain the basic structure of the original GreatWall. Then use linear regression to fit the generated 3D-model with higher resolution.
\item Expand the dimension of the brick images to include the structure-wise texture. Use GAN to generate the patterns of the missing part.
\end{rSubsection}

\end{rSection} 

\end{document}
